\chapter{系统性能实验}
\label{ch:exp}

\section{实验目的与环境}

本章在第二、三章给出的结构与参数设定基础上，对已实现的哈希系统进行全面的性能测试与功能有效性验证。本实验围绕操作延迟、k-kick 踢出次数、多层级 Cubby 向下合并（rebuild\_down）与向上拆分（rebuild\_up）三类指标展开，用以检验第三章各模块设计具体运行的效果\cite{bender2022otsh,bender2024sosa}， 通过实验数据与理论分析，多方面评估系统在不同数据规模、不同配置参数下的效果，为算法与系统设计实现的合理性提供数据支撑。

实验的系统代码基于 C++23 标准开发实现，元素规模 \(n\in\{10^3,10^4,10^5\}\)，Facility 规模 \(K\) 按分组取多档（基准 \(K=\log^3 n\)），k-kick 最大深度 \(k\in\{3,4,5\}\)（基准 \(k=4\)），NODE\_MAX\_BITS 常数 \(c=2\)。各组 workload(一次实验里对哈希表执行的操作序列及其构成方式) 在插入、命中查询与删除之间按预设比例混合执行，单次测量取各操作类型的平均延迟，以避免误差以及体现数据的代表性。

\section{实验方案}

\subsection{实验分组与工作量}

为了验证系统各模块的功能实现与效果，实验分为四组，如表~\ref{tab:exp-plan} 所示。规模 \(n\) 组与 \(K\) 组、\(k\) 组采用键均匀分布于各 Facility 的混合增删负载的形式；Cubby 动态重构组按纯插入、纯删除及不同插入/删除比例构造 workload，以分别触发 rebuild\_up 与 rebuild\_down。

\begin{table}[htbp]
  \centering
  \caption{实验分组与主要验证内容}
  \label{tab:exp-plan}
  \small
  \begin{tabular}{p{1.6cm}p{3.6cm}p{5.8cm}}
    \toprule
    分组 & 负载方式 & 主要验证内容 \\
    \midrule
    规模 \(n\) & 均匀分布、混合增删 & 操作延迟、踢出、rebuild \\
    Facility \(K\) & 均匀分布、混合增删 & 操作延迟、踢出 \\
    k-kick & 均匀分布、高负载 & 操作延迟、踢出 \\
    Cubby & Insert/Delete/混合比例 & rebuild\_down、rebuild\_up、重构耗时 \\
    \bottomrule
  \end{tabular}
\end{table}

\subsection{评价指标}

为量化评估系统的基础性能，实验记录并分析以下指标：

\begin{itemize}
  \item 操作延迟：插入、命中查询、删除的平均延迟。该指标用于衡量整个哈希系统的基础操作的性能；
  \item k-kick：workload 内累计踢出次数，以及单次插入/删除搬移上界 insert\_moved\_max、delete\_moved\_max。该指标用于衡量 k-kick 核心模块的具体表现，判断模块的基础踢出功能是否实现；
  \item j-tier Cubby：rebuild\_down（向下合并）与 rebuild\_up（向上拆分）触发次数。该指标用于验证多层级的 Cubby 的合并与拆分功能是否正常实现；
\end{itemize}

\section{规模 \texorpdfstring{$n$}{n} 实验}

固定 \(K=\log^3 n\)、\(k=4\)、\(c=2\)，令 \(n\) 从 \(10^3\) 增至 \(10^5\)。表~\ref{tab:exp-n-fixed} 给出各档测量结果。

\begin{table}[htbp]
  \centering
  \caption{规模 \(n\) 实验结果（\(K=\log^3 n\)，\(k=4\)，键均匀分布）}
  \label{tab:exp-n-fixed}
  \footnotesize
  \begin{tabular}{lrrrrrr}
    \toprule
    $n$ &
    ins ($\mu$s) & qry ($\mu$s) & del ($\mu$s) &
    总踢出 & rebuild\_down & rebuild\_up \\
    \midrule
    $10^3$  & 28.5 & 2.8  & 28.6 & 62     & 3     & 173   \\
    $10^4$  & 29.4 & 3.4  & 42.1 & 1188   & 63    & 911   \\
    $10^5$  & 51.2 & 6.1  & 197.3 & 10333 & 11500 & 20503 \\
    \bottomrule
  \end{tabular}
\end{table}

\noindent 表中 rebuild\_down 对应向下合并次数，rebuild\_up 对应向上拆分次数。

随 \(n\) 增大，命中查询延迟由 2.8\,$\mu$s 升至 6.1\,$\mu$s，相对 \(n\) 的三个数量级跨度增幅平缓，与 Local Query Router 的常数步定位一致\cite{bender2022otsh,morrison1968}。插入延迟由 28.5\,$\mu$s 增至 51.2\,$\mu$s；删除延迟由 28.6\,$\mu$s 增至 197.3\,$\mu$s，在高规模混合增删下删除路径承担更多 k-kick 与 tier 维护。workload 内累计踢出次数由 62 增至 \(1.0\times 10^4\)，rebuild\_down 由 3 增至 11500 次、rebuild\_up 由 173 增至 20503 次，说明随元素规模扩大，j-tiered Cubby 的向上拆分与向下合并均更频繁参与维护\cite{bender2024iceberg,pandey2023iceberght}。

\section{Facility 规模 \texorpdfstring{$K$}{K} 实验}

固定 \(n=10^4\)、\(k=4\)，比较 \(K\in\{64,256,1024\}\) 三档。表~\ref{tab:exp-K-fixed} 汇总三操作延迟与踢出指标。

\begin{table}[htbp]
  \centering
  \caption{Facility 规模 \(K\) 实验结果（\(n=10^4\)，键均匀分布）}
  \label{tab:exp-K-fixed}
  \small
  \begin{tabular}{lrrrrr}
    \toprule
    $K$ &
    ins ($\mu$s) & qry ($\mu$s) & del ($\mu$s) &
    总踢出 \\
    \midrule
    64    & 36.0 & 3.9 & 46.5 & 1841 \\
    256   & 34.3 & 3.3 & 33.6 & 1307 \\
    1024  & 33.4 & 3.3 & 36.7 & 1051 \\
    \bottomrule
  \end{tabular}
\end{table}

\(K\) 由 64 增至 1024 时，单 Facility 容量扩大，局部 k-kick 竞争减弱：累计踢出次数由 1841 降至 1051，插入延迟由 36.0\,$\mu$s 降至 33.4\,$\mu$s，删除延迟由 46.5\,$\mu$s 降至 33.6\,$\mu$s（\(K=1024\) 档为 36.7\,$\mu$s，与 256 档接近）；查询延迟在三档间不超过 3.9\,$\mu$s。\(K=64\) 与 \(K=1024\) 在踢出次数与增删延迟上同向改善，说明在本实现中扩大 Facility 规模可缓解局部冲突；实际部署时需结合目标负载率与并发访问粒度综合选取 \(K\)\cite{bender2022otsh,bender2024sosa}。

\section{k-kick 深度 \texorpdfstring{$k$}{k} 实验}

固定 \(n=10^4\)、\(K=256\)，对 \(k\in\{3,4,5\}\) 各执行相同高负载混合增删 workload。表~\ref{tab:exp-k-fixed} 给出完整结果。

\begin{table}[htbp]
  \centering
  \caption{k-kick 深度 \(k\) 实验结果（\(n=10^4\)，\(K=256\)，高负载）}
  \label{tab:exp-k-fixed}
  \small
  \begin{tabular}{lrrrrr}
    \toprule
    $k$ &
    ins ($\mu$s) & qry ($\mu$s) & del ($\mu$s) &
    总踢出 \\
    \midrule
    3 & 35.3 & 3.3 & 36.0 & 303  \\
    4 & 35.0 & 3.4 & 34.4 & 1201 \\
    5 & 36.3 & 3.4 & 33.8 & 1779 \\
    \bottomrule
  \end{tabular}
\end{table}

\noindent insert\_moved\_max 与 delete\_moved\_max 各档均不超过 \(k\)。

累计踢出次数由 303 增至 1779，随 \(k\) 单调增加，表明更大踢出深度在相同 workload 下触发更多搬移\cite{bender2022otsh,pagh2004cuckoo}。插入延迟在 35.0--36.3\,$\mu$s 之间，\(k=4\) 与 \(k=5\) 略高于 \(k=3\)；删除延迟在 33.8--36.0\,$\mu$s 之间波动。查询延迟在三档间均为 3.3--3.4\,$\mu$s，kick 深度主要作用于增删热路径，查询仍经 Router 常数步定位\cite{bender2022otsh,bender2024sosa}。

\section{多层级 Cubby 动态重构实验}

固定 \(n=10^4\)、\(K=256\)、\(k=4\)，按表~\ref{tab:exp-tier-fixed} 四种 workload 执行操作序列，统计 rebuild 次数与单次重构平均耗时。

\begin{table}[htbp]
  \centering
  \caption{多层级 Cubby 动态重构实验（\(n=10^4\)，\(K=256\)）}
  \label{tab:exp-tier-fixed}
  \small
  \begin{tabular}{lrrr}
    \toprule
    Workload      & rebuild\_down & rebuild\_up & 平均重构耗时 (ms) \\
    \midrule
    Insert-only   & 0     & 382   & 1.63 \\
    Delete-only   & 389   & 0     & 2.78 \\
    80/20 混合    & 8     & 381   & 1.60 \\
    50/50 混合    & 93    & 388   & 1.57 \\
    \bottomrule
  \end{tabular}
\end{table}

纯插入 workload 下 rebuild\_up（382 次）显著多于 rebuild\_down（0 次），说明 Cubby 数随元素增长时主要触发向上拆分；纯删除 workload 则相反，rebuild\_down（389 次）远多于 rebuild\_up（0 次），符合删除使 j-tier Cubby 稀疏后触发向下合并的路径\cite{bender2024iceberg,pandey2023iceberght}。80/20 混合负载以 rebuild\_up 为主（381 次），50/50 混合负载下两方向 rebuild 均非零且次数接近（93/388），随删除比例上升 rebuild\_down 占比逐步提高。四种 workload 的平均重构耗时在 1.57--2.78\,ms 之间，相对单次增删操作仍为次要开销，与算法~\ref{alg:tier-schedule} 仅在阈值 crossing 时批量导出重插的设计一致。

\section{本章小结}

本章节对系统性能进行了实验验证，包括规模 \(n\)、Facility 规模 \(K\)、k-kick 深度 \(k\)、多层级 Cubby 动态重构四个方面，验证了主要模块的正常运行效果，并给出了实验结果的分析。实验结果表明，设计的系统基本达到了预期目标，能够正常运行，在保证运行稳定的情况下同时具备良好的动态扩容能力。